% The complete finite-approximation, martingale and experiment comparison proofs.

\begin{annexitem}{annex:p17.compact-convex-tools}{thm:p17-compact-convex-tools}{Compact separation, affine minorants and integral Jensen}
\begin{proof}
Fix \(\epsilon>0\) and \(x\in C\). Choose
\(H>1+\max_C\phi\). The set
\[
                   E=\{(y,t):y\in C,\ \phi(y)\leq t\leq H\}
\]
is nonempty, compact and convex. The point
\(b=(x,\phi(x)-\epsilon)\) is outside it.
The proved finite-dimensional separation theorem gives
\((a,c)\) and a real \(\gamma\) such that
\[
             a\cdot y+ct\leq\gamma
                    <a\cdot x+c(\phi(x)-\epsilon)
                                      \quad((y,t)\in E).
\]
Applying the first inequality at \((x,\phi(x))\) shows
\(-c\epsilon>0\), so \(c<0\).
Consequently the globally affine function
\(\ell_x(y)=(a\cdot y-\gamma)/(-c)\) satisfies
\(\ell_x(y)\leq\phi(y)\) on \(C\) and
\(\ell_x(x)>\phi(x)-\epsilon\).
Continuity makes \(\phi-\ell_x<\epsilon\) on a relative neighborhood
of \(x\). Compactness supplies finitely many such neighborhoods
covering \(C\). Their corresponding maximum
\(\psi=\max_j\ell_{x_j}\) satisfies
\[
                           \phi-\epsilon<\psi\leq\phi
                                                   \quad\text{on }C.
\]
This proves uniform approximation, including lower-dimensional \(C\);
no interior-point hypothesis was used.

For the barycenter and Jensen assertions take a finite \(1/k\)-net
of points in \(C\) and send each \(y\in C\) to the first net point
within distance \(1/k\). The resulting finite-valued map \(r_k\)
is Borel measurable, takes values in \(C\), and obeys
\(\|r_k(y)-y\|_2\leq1/k\).
The vector \(\int r_k\,d\lambda\) is a finite convex combination
and hence belongs to \(C\). Its limit is \(\int y\,d\lambda\),
which belongs to \(C\) because \(C\) is closed.
Finite Jensen, obtained by induction on the number of positive weights,
gives
\[
        \phi\left(\int r_k\,d\lambda\right)\leq\int\phi(r_k)\,d\lambda.
\]
Uniform continuity on \(C\) passes both sides to the asserted limit.

Finally let the two finite laws be supported by a finite set \(D\),
and put \(C=\operatorname{conv}(D)\), a compact convex set.
Each continuous convex test on \(C\) is uniformly approximated there
by the global convex functions just constructed. The global test
inequality therefore passes to every such test on \(C\).
All hypotheses of the finite transport theorem now hold.
\end{proof}
\end{annexitem}

\begin{annexitem}{annex:p17.continuous-law-identification}{thm:p17-continuous-law-identification}{Closed-set approximation and uniform continuity}
\begin{proof}
For a nonempty closed \(F\subseteq\mathbb R^d\), the function
\(\operatorname{dist}(x,F)=\inf_{z\in F}\|x-z\|_2\)
is 1-Lipschitz by the triangle inequality, and vanishes exactly on \(F\).
Thus
\[
              f_k(x)=\max(0,1-k\operatorname{dist}(x,F))
                               \downarrow\mathbf1_F(x)
\]
is a bounded continuous approximation.
Equality of integrals and dominated convergence give equality of the
two probabilities on \(F\). The empty closed set has mass zero.
Closed sets form a pi-system, contain the whole space, and generate
the Borel sigma-algebra by complementation. The measure-uniqueness
theorem proves equality of the probabilities.
On a compact subspace \(K\), use the same distance construction for
its relatively closed subsets, which are closed in the ambient space.
This proves the second assertion.

A continuous function on a compact set is uniformly continuous:
otherwise there would be \(\epsilon>0\) and \(x_j,y_j\in K\) with
\(\|x_j-y_j\|_2\to0\) but \(|f(x_j)-f(y_j)|\geq\epsilon\).
A convergent subsequence of \(x_j\), and then the same limit for \(y_j\),
contradicts continuity. Hence \(\omega_{f,K}(r)\to0\) as \(r\downarrow0\).
Pointwise moves of length at most \(r\) change the function by at most
this modulus; integrating the bound gives the stated control.
\end{proof}
\end{annexitem}

\begin{annexitem}{annex:p17.centroid-refinement}{thm:p17-centroid-refinement}{Conditional cell means and every refinement row}
\begin{proof}
Realize a vector \(Z\) with law \(\mu\) on the bounded Borel set \(X\).
Only finitely many level-\(N\) cells can have positive mass, since
each intersects one fixed bounded box containing \(X\).
Define \(C_N(Z)=c_{N,a}\) in each positive-mass cell; on all other
cells choose, for definiteness, the zero vector.
The exceptional union has \(\mu\)-mass zero. Each coordinate of a
positive-mass centroid lies in the closure of its cell, since its
coordinate is the average of numbers in that interval.
Also it lies in a fixed closed box containing \(X\), by the
barycenter result for that box. The law of \(C_N(Z)\) is \(\mu_N\).

The sigma-algebra of level-\(N\) cells is contained in that at level
\(N+1\), because each cell is a disjoint union of its \(2^d\)
children. Conditional cell averaging and the tower identity give
\[
 E(Z\mid\text{level-}N\text{ cells})=C_N(Z),\qquad
 E(C_{N+1}(Z)\mid\text{level-}N\text{ cells})=C_N(Z).
\]
For \(a\in A_N\), define the finite kernel
\[
 T_N(a,b)=
   \frac{P(C_N(Z)=a,\ C_{N+1}(Z)=b)}{P(C_N(Z)=a)}
                                               \quad(b\in A_{N+1}).
\]
Its denominator is positive by the definition of \(A_N\), its row
sums to one, and its marginal identity is \(\mu_NT_N=\mu_{N+1}\).
Conditioning the preceding equality further on \(C_N(Z)\) gives
\(\sum_bT_N(a,b)b=a\). This handles coincidences of centroids from
different cells without choosing a parent.
Jensen for each of these finite rows yields the first convex-order
inequality. Jensen for the conditional probability of \(Z\) in each
positive cell yields \(\mu_N\preceq_{\mathrm{cx}}\mu\), and the
same argument at level \(N+1\) gives the displayed chain.

Almost surely \(Z\) and \(C_N(Z)\) are in the closure of one cell,
so their Euclidean distance is at most \(\sqrt d\,h_N\).
Uniform continuity on the common compact box gives
\[
 |\!E f(C_N(Z))-Ef(Z)\!|
                  \leq\omega_{f,K}(\sqrt d\,h_N).
\]
All integrals here are finite; the bound proves the final assertion.
\end{proof}
\end{annexitem}

\begin{annexitem}{annex:p17.vertex-refinement}{thm:p17-vertex-refinement}{Normalized splitting and the nested-grid identity}
\begin{proof}
Within a cell, \(0\leq\lambda_{N,i}(t)<1\).
Summing the product weights over \(\epsilon\in\{0,1\}^d\) gives
\(\prod_i[(1-\lambda_{N,i})+\lambda_{N,i}]=1\).
They are nonnegative. The mean of coordinate \(i\) of the chosen
vertex is
\[
      h_Na_i(1-\lambda_{N,i})
        +h_N(a_i+1)\lambda_{N,i}=t_i;
\]
the other coordinate weights sum to one. Thus the kernel is mean
preserving at every point. Each cell and each weight restricted to it
is Borel; a countable sum over cells proves measurability for every
Borel target event.

To prove \(U_{N+1}U_N=U_N\), first fix a coarse cell closure and a
coarse vertex \(v\). On this closure the probability \(U_N(t,\{v\})\)
is the product of the appropriate affine coordinate weights; for
other vertices it is zero. The same expression remains valid on a
cell face. Indeed, there the weight on an incompatible side is zero,
and both adjacent cell formulas place the mass on the common face.
Hence the half-open cell convention does not change the measure
on a boundary point.

A finer cell lies inside one coarse cell closure. Splitting \(t\)
onto its fine vertices makes their coordinates independent with
means \(t_i\). The expectation of a product of affine functions of
the separate coordinates is therefore the product of their
expectations, equal to the coarse product weight at \(t\).
This proves the composition identity on each coarse vertex.
Both composed and direct probabilities have support in the same
finite set of coarse vertices, so it proves equality as measures.

Consequently
\[
                   \nu_NU_{N-1}
                     =\nu U_NU_{N-1}
                     =\nu U_{N-1}=\nu_{N-1}\quad(N\geq2).
\]
If \(b\in B_N\) and a coarse vertex \(c\) has
\(U_{N-1}(b,\{c\})>0\), then
\(\nu_{N-1}(\{c\})\geq\nu_N(\{b\})U_{N-1}(b,\{c\})>0\).
Thus each restricted row really takes values in \(B_{N-1}\).
The marginal and mean identities give the finite \(V_N\).
Jensen first for \(U_{N+1}\) and then for \(V_{N+1}\) proves
\(\nu\preceq_{\mathrm{cx}}\nu_{N+1}\preceq_{\mathrm{cx}}\nu_N\).

Every chosen vertex is within \(\sqrt d\,h_N\) of its input.
If \(X\subseteq[-R,R]^d\), all such vertices for \(N\geq1\)
belong, for example, to \(K=[-R-1,R+1]^d\).
Integrating the uniform-continuity bound over each splitting row
and then over \(\nu\) proves the stated approximation inequality.
\end{proof}
\end{annexitem}

\begin{annexitem}{annex:p17.compatible-transports}{thm:p17-compatible-transports}{The full diagonal construction of the source's compatible kernels}
\begin{proof}
The previous chains and the hypothesis imply
\(\mu_D\preceq_{\mathrm{cx}}\nu_D\) for every \(D\geq1\).
Both laws have finite support in \(A_D\cup B_D\).
The compact-convex lemma and the finite transport theorem supply a
mean-preserving finite comparison. Restrict its rows to \(A_D\).
No row of positive \(\mu_D\)-mass can put positive mass outside \(B_D\),
since the target law has zero mass there and all summands are
nonnegative. Thus obtain a kernel \(J_D:A_D\to B_D\) with
\(\mu_DJ_D=\nu_D\).

For \(D>N\), form
\[
 H_{N,D}=T_NT_{N+1}\cdots T_{D-1}\,
               J_D\,V_DV_{D-1}\cdots V_{N+1}.
\]
It is a row-stochastic matrix of the fixed shape \(A_N\times B_N\).
Composition preserves the marginal law and the row barycenter:
if two consecutive rows have mean equal to their respective inputs,
finite summation gives the same property for their product.
It follows that
\[
             \mu_NH_{N,D}=\nu_N,\qquad
             \sum_bH_{N,D}(a,b)b=a
                                         \quad(a\in A_N).
\]
For \(D>N+1\), associativity also gives the exact relation
\[
                         H_{N,D}=T_NH_{N+1,D}V_{N+1}.
\]

For fixed \(N\), these matrices lie in the finite-dimensional compact
cube \([0,1]^{A_N\times B_N}\).
Select an infinite subsequence of the positive integers \(D\)
along which \(H_{1,D}\) converges; from that choose one along
which \(H_{2,D}\) converges, and continue. Choose the diagonal
indices \(D_j\) strictly increasing and at least \(j+1\), with the
tail of the chosen sequence contained in each of the first \(j\)
subsequences. Then for every fixed \(N\),
\(H_{N,D_j}\to W_N\) entry by entry.
All row sums, nonnegativity constraints, marginal equations and
mean equations are finite linear conditions, so they pass to the limit.
The last displayed compatibility relation passes to the limit too,
because multiplication by fixed finite matrices is continuous.
This gives all asserted equations simultaneously.
\end{proof}
\end{annexitem}

\begin{annexitem}{annex:p17.blackwell-joint-process}{thm:p17-blackwell-joint-process}{Consistency, both filtrations and the terminal mean identity}
\begin{proof}
For each \(N\) give the ordered variables
\((L_1,\ldots,L_N,R_N,\ldots,R_1)\) the finite probability masses
\[
 \mu_1(l_1)
 \left[\prod_{i=1}^{N-1}T_i(l_i,l_{i+1})\right]
 W_N(l_N,r_N)
 \left[\prod_{i=2}^{N}V_i(r_i,r_{i-1})\right].
\]
For \(N=1\) both empty products equal one.
Successively summing rows shows total mass one.
The marginal laws are \(\mu_i,\nu_i\), by the established transport
identities. Reorder the coordinates as
\((L_1,R_1,\ldots,L_N,R_N)\).
Marginalizing the level-\((N+1)\) law over its two extra vectors
replaces the three central kernels
\(T_NW_{N+1}V_{N+1}\) by \(W_N\).
All other factors stay unchanged. This is exactly the level-\(N\)
law, proving consistency.
Flatten each vector into its finitely many real coordinates.
Kolmogorov extension
\cref{thm:p15-kolmogorov-countable} supplies the countable process
with all these finite distributions.

The \(L_i,R_i\) take values in their specified finite supports
almost surely; the countable intersection of these events has
probability one. All supports are in one compact box, so each
coordinate has a common absolute bound. The finite factorization
implies
\[
        E(L_{N+1}\mid L_1,\ldots,L_N)=L_N
\]
by the row-mean identity of \(T_N\). Thus \(L_N\) is a forward
martingale for \(\mathcal F_N=\sigma(L_1,\ldots,L_N)\).

For the reverse assertion fix \(N\). The finite factorization of
\((R_m,\ldots,R_{N+1},R_N)\), for every \(m>N\), gives
\[
 E\bigl((R_N-R_{N+1})\mathbf1_A\bigr)=0
                \quad\bigl(A\in\sigma(R_{N+1},\ldots,R_m)\bigr),
\]
coordinate by coordinate, using the mean identity of \(V_{N+1}\).
For one coordinate the class of events satisfying this identity
is a lambda-system: its total integral is zero, complements preserve
the identity, and countable disjoint unions do so by dominated
convergence with a bounded integrand. It contains the pi-system
\(\bigcup_{m>N}\sigma(R_{N+1},\ldots,R_m)\).
The pi-lambda theorem extends the identity to
\(\mathcal G_{N+1}=\sigma(R_{N+1},R_{N+2},\ldots)\).
Since \(R_{N+1}\) is measurable there, this proves
\(E(R_N\mid\mathcal G_{N+1})=R_{N+1}\).
The sigma-algebras \(\mathcal G_N\) decrease, so this is the
specified backward martingale.

The bounded forward martingale is uniformly integrable; the backward
martingale also satisfies its full convergence theorem.
The finite-coordinate versions of
\cref{thm:p15-closed-martingale,thm:p15-backward-martingale-convergence}
therefore give bounded limits \(L_*,R_*\), almost surely and in \(L^1\).
Choose measurable versions equal to the finite limits on their
convergence events, with a fixed box point elsewhere.
For every continuous \(f\) on the box, bounded convergence gives
\[
 E f(L_*)=\lim_NEf(L_N)=\lim_N\int f\,d\mu_N=\int f\,d\mu .
\]
The centroid estimate gives the last equality. The vertex estimate
gives the analogous equality for \(R_*\) and \(\nu\).
Continuous-test uniqueness identifies their laws.

Finally fix \(k\). For \(N\geq k\), composition of the row-mean
identities in the finite chain gives
\[
                           E(R_N\mid\mathcal F_k)=L_k.
\]
Passing to the \(L^1\) limit, using conditional contraction in each
coordinate, yields \(E(R_*\mid\mathcal F_k)=L_k\).
The upward conditional-limit theorem identifies the limit of these
conditional expectations as
\(E(R_*\mid\mathcal F_\infty)\), where
\(\mathcal F_\infty=\sigma(\bigcup_k\mathcal F_k)\).
Their already known limit is \(L_*\); hence
\(E(R_*\mid\mathcal F_\infty)=L_*\).
Since \(L_*\) is \(\mathcal F_\infty\)-measurable, conditioning further
on \(\sigma(L_*)\) proves \(E(R_*\mid L_*)=L_*\).
\end{proof}
\end{annexitem}

\begin{annexitem}{annex:p17.general-convex-transport}{thm:p17-general-convex-transport}{The everywhere defined barycentric kernel}
\begin{proof}
From convex order the preceding construction gives a bounded pair
with laws \(\mu,\nu\) and \(E(R_*\mid L_*)=L_*\).
Both variables belong to the Borel set \(X\) almost surely.
Choose \(x_0\in X\), and replace either variable by \(x_0\) on its
measurable exceptional nonmembership set.
They are now \(X\)-valued everywhere and all distributional and
conditional identities are unchanged.

As a Borel subset of Euclidean space, \(X\) is standard Borel by
\cref{thm:p15-metric-borel-representation}.
The value-parametrized conditional-law theorem
\cref{thm:p15-value-conditional-law} gives a Borel probability kernel
\(K_0(x,dy)\) for \(R_*\) given \(L_*=x\), valid as a probability
measure at every \(x\in X\), and with \(\mu K_0=\nu\).
Each coordinate of
\[
                              b(x)=\int_Xy\,K_0(x,dy)
\]
is bounded and Borel measurable by kernel integration.
The conditional-mean identity says \(b(x)=x\) for \(\mu\)-almost
every \(x\). The set \(D=\{x\in X:b(x)=x\}\) is Borel and
\(\mu(D)=1\). Define
\[
 K(x,A)=
 \begin{cases}
 K_0(x,A),&x\in D,\\
 \mathbf1_A(x),&x\notin D.
 \end{cases}
\]
For every Borel \(A\subseteq X\) this is measurable, and each row
is a probability measure. Its mean is \(x\) on both pieces.
Changing a kernel on the \(\mu\)-null complement of \(D\) preserves
\(\mu K=\nu\). This proves the full every-input assertion.

Conversely, given such a kernel, the joint-measure theorem supplies
\(\gamma(dx,dy)=\mu(dx)K(x,dy)\), whose marginals are \(\mu,\nu\).
Its coordinate conditional expectation is \(E(R\mid L)=L\), by
the kernel mean equation and event integration.
For a global continuous convex \(\phi\), all measures lie in a
fixed compact convex box. Jensen for each row gives
\(\phi(x)\leq\int\phi(y)K(x,dy)\).
Integrating over \(\mu\) yields convex order.
The same argument starts with an arbitrary joint law and its
conditional kernel, proving the joint-law equivalence.

If the premise is instead the source's inequality for every continuous
convex function on \(\operatorname{conv}(X)\), restrict each global
test to that convex hull. The resulting inequality is exactly the
sufficient global-test premise already proved.
Thus the source's bounded-set implication is obtained without a
closedness assumption on \(X\) and without assigning a nonexistent
barycenter to a zero-mass cell.
\end{proof}
\end{annexitem}

\begin{annexitem}{annex:p17.transport-calculations}{thm:p17-transport-calculations}{The complete probability and boundary calculations}
\begin{proof}
At \(N=1\), \(h_N=1/2\) and the indicated point is in the cell
\([0,1/2)^2\). Its two barycentric coordinates are \(1/2\) and \(3/4\).
Their product weights in the stated vertex order are
\((1/2)(1/4),(1/2)(1/4),(1/2)(3/4),(1/2)(3/4)\).
They sum to one. The first-coordinate mean is
\((1/2)(1/8+3/8)=1/4\), and the second is
\((1/2)(3/8+3/8)=3/8\).
Using unnormalized distances in place of barycentric coordinates
would multiply the total mass by \(h_N^2\) and would not define
a probability kernel.

For the second example the first cell \([0,h_N)\) has mass
\((1+h_N)/2\), and its first moment is \(h_N^2/4\).
Its centroid is therefore \(h_N^2/[2(1+h_N)]>0\).
Every other positive-mass cell has a strictly positive centroid too.
Thus every atom of \(\mu_N\) lies in \((0,\infty)\), giving mass one,
whereas precisely the continuous half of \(\mu\) lies in that open set.
This does not contradict the proved continuous-test convergence.

In the third example, the kernel weights are nonnegative and sum to one,
and the row mean is \(0(1-x)+1x=x\).
Under the uniform input law each endpoint receives integrated mass
\(1/2\). The second moments of the input and output are \(1/3\)
and \(1/2\), respectively. A mean-preserving kernel in the reverse
direction would imply \(1/2\leq1/3\) by Jensen for \(x^2\),
which is false.

For the last example, only the row at input zero affects the joint
law of two variables equal to zero almost surely. The constant
kernel \(\delta_0\) thus represents that conditional law, but has
mean zero at every \(x>0\). The kernel \(\delta_x\) is measurable
since \(\delta_x(A)=\mathbf1_A(x)\), has row mean \(x\), and
still sends \(\delta_0\) to \(\delta_0\). This verifies the distinction.
\end{proof}
\end{annexitem}

\begin{annexitem}{annex:p17.general-standard-measure}{thm:p17-general-standard-measure}{Densities, policies and canonical observations}
\begin{proof}
For every state \(s\), \(\rho(A)\geq\pi_sP_s(A)\);
hence \(P_s\ll\rho\). Radon--Nikodym gives a nonnegative density
\(r_s=dP_s/d\rho\), integrable with integral one, and finite
\(\rho\)-almost everywhere. The finite sum
\(\sum_s\pi_sr_s\) integrates on every event to \(\rho\) and
therefore equals one almost everywhere by density uniqueness.
Outside the common measurable null set where these properties fail,
put \(p_s=\pi_sr_s\); on that set choose a fixed simplex point.
Every \(P_s\) gives that set mass zero as well.
Thus all statewise integral identities are preserved.
The law \(\mu\) of \(p\) under \(\rho\) satisfies
\[
                             \int q_s\,d\mu(q)=\int p_s\,d\rho=\pi_s.
\]

For a finite action set enumerate its elements once. Transforming
the policy integral by the densities gives
\[
 J(\pi,\mathsf E,u,\delta)
       =\int_Z\sum_d\delta(d\mid z)
                       \sum_sp_s(z)u(d,s)\,d\rho(z).
\]
The inner sum is a convex combination of the finitely many displayed
action values, so it is at most their maximum.
Choose the first action attaining that maximum. Its selection sets
are finite intersections of measurable comparison sets, so this
defines a measurable deterministic policy attaining equality.
Transport through the law of \(p\) proves the value formula.

Conversely, for a simplex law \(\mu\) with barycenter \(\pi\),
each measure \(P_s^\mu(A)=\int_Aq_s/\pi_s\,d\mu(q)\)
is nonnegative and has total mass one. Their \(\pi\)-mixture is
\(\mu\) because \(\sum_sq_s=1\). Its posterior coordinates are
the coordinate functions \(q_s\), so its posterior law is \(\mu\).

It remains to compare every compact convex loss set, including when
the original observation space is arbitrary measurable.
Let \(f:Z\to C\) be any measurable decision.
Under \(\rho\), the vector \(f\) takes values in the standard Borel
compact set \(C\), so the value-parametrized conditional-law theorem
gives a probability kernel \(J(q,da)\) for \(f\) given \(p=q\),
at every \(q\in\Delta(S)\).
Its barycenter \(g(q)=\int_Ca\,J(q,da)\) is Borel measurable,
bounded, and belongs to \(C\) by the compact-convex lemma.
It is a version of \(E_\rho(f\mid p)\). Since \(p_s\) is bounded
and known given \(p\), the conditional pullout identity gives
\[
 \begin{aligned}
 \int_Z f_s\,dP_s
   &=\frac1{\pi_s}\int_Zp_sf_s\,d\rho
    =\frac1{\pi_s}\int_Zp_sg_s(p)\,d\rho\\
   &=\int_{\Delta(S)}g_s(q)\,\frac{q_s}{\pi_s}\,d\mu(q).
 \end{aligned}
\]
This is the canonical experiment's loss for \(g\).
Conversely every canonical decision \(g\) gives the original
decision \(g\circ p\) and the same displayed identity.
The loss sets are therefore equal.

For the uniform prior, \(\rho=(\sum_jP_j)/n\).
The Radon--Nikodym chain identity yields
\[
        p_s=\frac1n\frac{dP_s}{d\rho}
                         =\frac{dP_s}{d(\sum_jP_j)}
                                  \quad\text{almost everywhere}.
\]
The law and barycenter are exactly those in the source.
\end{proof}
\end{annexitem}

\begin{annexitem}{annex:p17.compact-measure-selection}{thm:p17-compact-measure-selection}{Consistent limiting cell masses and a compactly supported limit}
\begin{proof}
At each dyadic level \(N\geq1\), intersect the half-open cubes of
\(\mathbb R^m\) with \(K\). The nonempty intersections form a finite
partition, indexed by a finite set \(I_N\subseteq\mathbb Z^m\).
Each cell at the next level has a unique parent; each parent has
at most \(2^m\) children. Choose a representative \(z_{N,a}\in K\)
in each nonempty cell.
If the given measures are \(\lambda_j\), write their cell masses as
\(p_{j,N}(a)\).
These finite vectors lie in compact cubes.
Successive subsequence selection and then the diagonal selection
give indices \(j_l\uparrow\infty\) such that
\[
                      p_{j_l,N}(a)\longrightarrow p_N(a)
                                      \quad\text{for every }N,a.
\]
The limits are nonnegative, sum to one at each level, and satisfy
\[
                   p_N(a)=\sum_{b:\,\operatorname{parent}(b)=a}p_{N+1}(b),
\]
since each equality before the limit is a finite sum.

For the finite sequence of cell indices \(a_1,\ldots,a_N\),
assign mass \(p_N(a_N)\) if all parent relations hold, and zero otherwise.
The unique ancestors of a cell show total mass one.
Summing over the last index uses the preceding parent identity and
gives precisely the shorter law.
View the integer-vector indices as finitely supported real vectors;
Kolmogorov extension gives their countable joint realization.
Every parent relation holds almost surely, and their countable
intersection has probability one.

On that event, for \(M\geq N\), the representatives
\(z_{M,a_M}\) and \(z_{N,a_N}\) lie in the same level-\(N\) cell,
so their distance is at most \(\sqrt m\,2^{-N}\).
They form a Cauchy sequence. Its limit \(Z\) belongs to \(K\),
since all representatives do and \(K\) is closed.
Define \(Z\) to equal a fixed point of \(K\) on the exceptional set.
It is a measurable limit version and has a Borel law \(\lambda\) on \(K\).

For \(f\) continuous on \(K\),
\[
 \left|\int f\,d\lambda_j-\sum_{a\in I_N}p_{j,N}(a)f(z_{N,a})\right|
                         \leq\omega_{f,K}(\sqrt m\,2^{-N}).
\]
The same bound holds between \(Ef(Z)\) and
\(\sum_ap_N(a)f(z_{N,a})\), because the limit remains in the
closure of every cell in its path.
For fixed \(N\) the finite sums converge along \(j_l\).
The two bounds yield
\[
 \limsup_l\left|\int f\,d\lambda_{j_l}-\int f\,d\lambda\right|
                    \leq2\omega_{f,K}(\sqrt m\,2^{-N}).
\]
Let \(N\to\infty\).
The selected subsequence depends only on the countable cell masses,
so this argument applies to every continuous \(f\) on that same
subsequence, as required.
\end{proof}
\end{annexitem}

\begin{annexitem}{annex:p17.general-risk-compactness}{thm:p17-general-risk-compactness}{Passing a sequence of decisions through joint probability laws}
\begin{proof}
Use the canonical representation with a fixed full-support prior
\(\pi\) and posterior law \(\mu\).
Nonemptiness follows by taking a constant decision in \(C\).
Every coordinate of an attainable loss vector is bounded by
\(\max_{a\in C}|a_s|\), since \(P_s\) is a probability.
If \(f,g\) are decisions and \(0\leq t\leq1\),
the pointwise decision \(tf+(1-t)g\) is measurable and lies in \(C\).
Linearity of all integrals proves convexity.

Suppose a sequence of attainable vectors \(v_j\) tends to \(v\).
Choose canonical decisions \(g_j:\Delta(S)\to C\) attaining them
and form the joint laws
\[
                     \gamma_j=\operatorname{law}_{q\sim\mu}(q,g_j(q))
\]
on the compact set \(K=\Delta(S)\times C\).
The compact-selection lemma gives a subsequence converging under
every continuous test to a probability \(\gamma\) on \(K\).
For continuous \(h\) on \(\Delta(S)\),
\(\int h(q)\,d\gamma_j(q,a)=\int h\,d\mu\).
Passing to the limit and using continuous-test uniqueness shows
that the first marginal of \(\gamma\) is \(\mu\).
For each state the function \(q_sa_s/\pi_s\) is continuous on \(K\);
therefore
\[
                        v_s=\int_K\frac{q_sa_s}{\pi_s}\,d\gamma(q,a).
\]

Take a regular conditional law \(J(q,da)\) of the second coordinate
given the first. It is a probability on \(C\) at every \(q\).
Its barycenter \(g(q)\) is Borel and belongs to \(C\).
Kernel integration now gives
\[
                         v_s=\int_{\Delta(S)}
                                   \frac{q_sg_s(q)}{\pi_s}\,d\mu(q).
\]
Thus \(g\) is a canonical decision attaining \(v\).
This proves closedness. A closed bounded subset of finite-dimensional
Euclidean space is compact by the compactness result already supplied.
Returning through the canonical representation proves the assertion
for the original observation space.
\end{proof}
\end{annexitem}

\begin{annexitem}{annex:p17.general-decision-comparison}{thm:p17-general-decision-comparison}{The full equivalence and uniqueness of the standard measure}
\begin{proof}
Write \(\mu=\mu_{\mathsf E,\pi}\) and \(\nu=\mu_{\mathsf F,\pi}\)
for one fixed full-support reference prior.
First, loss-set inclusion implies the value comparison at every
prior \(\lambda\), including those with zero entries.
For a finite payoff menu take \(C\) to be the convex hull of its
negative payoff vectors \((-u(d,s))_{s=1}^n\).
The optimal value is
\[
                         \max_{v\in B(\mathsf E,C)}-\sum_s\lambda_sv_s.
\]
To verify this equality, represent all state measures using the
full-support reference mixture and its posterior \(p\).
For any \(C\)-valued decision, its pointwise payoff under the weights
\(\lambda_sp_s/\pi_s\) is at most the maximum of the original finite
menu. The first maximizing menu element is a measurable decision and
attains equality, as in the preceding value proof.
Thus the displayed optimization is exactly the finite-policy value.
Inclusion of the loss sets proves condition 4 from condition 1.
Condition 4 immediately gives condition 3.

Under condition 3, every maximum of finitely many linear payoff
functions on \(\Delta(S)\) has at least as large an integral under
\(\mu\) as under \(\nu\).
An affine function \(a\cdot q+b\) on the simplex is linear there
with coefficients \(a_s+b\), since \(\sum_sq_s=1\).
The compact-convex approximation lemma therefore extends the inequality
to every continuous convex function on \(\Delta(S)\), and in particular
to every restricted global convex test. This is condition 2.

Assume condition 2. The bounded transport theorem on
\(X=\Delta(S)\) gives a mean-preserving kernel \(T\) such that
\(\nu T=\mu\). Form the joint law
\(\gamma(dq,dp)=\nu(dq)T(q,dp)\), and take its reverse conditional
kernel \(R(p,dq)\); both coordinates lie in the standard Borel simplex.
Its first-coordinate marginal in the \(p\) direction is \(\mu\),
and, for any Borel \(A\subseteq\Delta(S)\),
\[
 \begin{aligned}
 \int_{\Delta(S)}\frac{p_s}{\pi_s}R(p,A)\,d\mu(p)
   &=\frac1{\pi_s}\int_{q\in A}\left(\int p_sT(q,dp)\right)d\nu(q)\\
   &=\int_A\frac{q_s}{\pi_s}\,d\nu(q).
 \end{aligned}
\]
Thus \(R\) sends every state measure of the more informative
canonical experiment to the corresponding state measure of the less
informative one. For a canonical loss decision \(g(q)\in C\), put
\(h(p)=\int_{\Delta(S)}g(q)R(p,dq)\), coordinate by coordinate.
This bounded Borel vector belongs to \(C\) by
the barycenter lemma applied to the law of \(g(q)\).
The preceding statewise identity shows that its risk under the
first canonical experiment equals the risk of \(g\) under the second.
Canonical representation transfers this inclusion to the original
experiments. This proves condition 1 and closes the cycle.
Applying the same cycle at any other full-support prior proves the
stated prior independence.

If two standard measures are equal, their canonical representations
give equal loss sets. Conversely equal loss sets give convex domination
in both directions at the uniform prior, hence equal integrals of every
global continuous convex test for the two simplex laws.
For \(t\in\mathbb R^n\), define the nonnegative convex function
\[
                       D_t(q)=\max(0,q_1-t_1,\ldots,q_n-t_n).
\]
Both \(D_t\) and \((D_t-\epsilon)^+\) are global continuous convex
functions. Therefore the integrals of
\[
       \frac{D_t-(D_t-\epsilon)^+}{\epsilon}
                            =\min(D_t/\epsilon,1)
\]
agree for the two laws. As \(\epsilon\downarrow0\), this expression
tends to \(\mathbf1_{\{D_t>0\}}\), bounded by one.
Their probabilities of
\(\{q:q_s\leq t_s\text{ for all }s\}\) consequently agree.
Lower coordinate rectangles, with the whole space adjoined, form a
generating pi-system for the Borel sets of \(\mathbb R^n\).
The uniqueness theorem identifies the two standard measures.
\end{proof}
\end{annexitem}

\begin{annexitem}{annex:p17.general-garbling}{thm:p17-general-garbling}{Recovering observations and reversing the posterior transformation}
\begin{proof}
If \(P_sG=Q_s\), simulate a compact-convex loss decision \(f(y)\)
by \(h(z)=\int f(y)G(z,dy)\).
Kernel integration makes \(h\) measurable, and its value belongs
to \(C\) by the barycenter lemma.
Integrating under \(P_s\) gives exactly \(\int f_s\,dQ_s\).
Thus the general decision-comparison theorem applies.

The source's opposite direction between the two posterior laws can
also be seen directly. Let the reference prior be \(\pi\), the
mixtures be \(\rho,\eta\), and the two posterior maps be \(p(z),q(y)\).
Under the joint observation law \(\rho(dz)G(z,dy)\),
the second marginal is \(\eta\).
For any measurable \(A\subseteq Y\),
\[
 \begin{aligned}
 E\bigl(p_s(Z)\mathbf1_{\{Y\in A\}}\bigr)
   &=\int_Zp_s(z)G(z,A)\,d\rho(z)\\
   &=\pi_sQ_s(A)=\int_Aq_s(y)\,d\eta(y).
 \end{aligned}
\]
Hence \(E(p_s(Z)\mid Y)=q_s(Y)\), and the tower identity gives
\(E(p(Z)\mid q(Y))=q(Y)\).
A regular conditional law of \(p(Z)\) given \(q(Y)\) is a
kernel from the less informative posterior to the more informative
posterior, transporting \(\nu\) to \(\mu\).
Its mean equals the input \(\nu\)-almost everywhere.
The measurable replacement by \(\delta_q\) on the exceptional set,
as in the transport theorem, makes it mean preserving everywhere.

Conversely suppose the equivalent posterior comparison holds.
Let \(T:\Delta(S)\to\Delta(S)\) be the mean-preserving kernel
from \(\nu\) to \(\mu\), and let \(R(p,dq)\) be the reverse
conditional kernel of the joint law \(\nu(dq)T(q,dp)\).
The preceding proof gives the statewise canonical identity
\[
                 \left(\frac{p_s}{\pi_s}\mu(dp)\right)R
                                =\frac{q_s}{\pi_s}\nu(dq).
\]
Under the less informative observation mixture \(\eta\), take a
regular conditional law \(J(q,dy)\) of the actual observation \(Y\)
given its posterior value \(q(Y)\).
This exists with rows on \(Y\) at every \(q\) when \(Y\) is
standard Borel. Kernel integration and the conditional pullout
identity give, for every \(A\in\mathcal Y\),
\[
 \int_{\Delta(S)}\frac{q_s}{\pi_s}J(q,A)\,d\nu(q)
       =\int_Y\frac{q_s(y)}{\pi_s}\mathbf1_A(y)\,d\eta(y)=Q_s(A).
\]
Define the actual simulation by
\[
                          G(z,A)=\int_{\Delta(S)}
                                         J(q,A)R(p(z),dq).
\]
Composition of probability kernels is a probability kernel:
measurability follows by integrating the bounded measurable
function \(q\mapsto J(q,A)\), total mass is one, and countable
additivity follows from monotone convergence for disjoint target
sets. Composition with the measurable map \(p\) preserves measurability.
Using the density identity for the law of \(p\) under each \(P_s\)
and then the two displayed statewise identities proves \(P_sG=Q_s\).
All null-input rows are already probability measures.

This proof requires an observation-recovery kernel on the target
space. It applies, more generally, whenever that kernel is given.
The loss-set argument did not need this hypothesis because it
conditioned a finite-dimensional action vector, whose target is
always standard Borel. The distinction is therefore an explicit
measurable-space condition on simulation.
\end{proof}
\end{annexitem}

\begin{annexitem}{annex:p17.continuous-experiment-example}{thm:p17-continuous-experiment-example}{Every density, optimal value and simulation}
\begin{proof}
Both proposed densities are nonnegative and integrate to one:
\(\int_0^12x\,dx=1=\int_0^12(1-x)\,dx\).
At prior \((1/2,1/2)\), their mixture density is
\(x+(1-x)=1\). The posterior coordinates are therefore \(x,1-x\).
For unit payoff for a correct state guess and zero otherwise,
choose state 1 for \(x\geq1/2\) and state 2 for \(x<1/2\).
The value is
\[
                  \int_0^1\max(x,1-x)\,dx
                       =2\int_{1/2}^1x\,dx=\frac34.
\]
The no-information signal leaves the two posterior probabilities
equal to \(1/2\), so its optimal value is \(1/2\).
The constant map to that signal is the explicit garbling.

For the perfectly revealing experiment with output \(s\) in state \(s\),
the kernel assigning law \(P_s\) to the observed label \(s\)
sends each point mass \(\delta_s\) to \(P_s\).
It is one kernel depending only on the observed label.

Any kernel from a one-point observation space has the same output
law in every state. The proposed two densities do not:
\[
             P_1([1/2,1])=\frac34,\qquad
             P_2([1/2,1])=\frac14.
\]
Therefore no reverse simulation exists.
Finally, at prior \((1,0)\), every experiment has value
\(\max_du(d,1)\) in every finite payoff problem: the first state's
payoff can never exceed this number, and the maximizing constant
action attains it. Thus all value comparisons at that prior hold
even though the just-proved simulation obstruction remains.
This proves the zero-prior limitation directly.
\end{proof}
\end{annexitem}
